\section{Introduction}\label{sec:intro}

A common feature-selection recipe in applied machine learning is to add columns that are known to be useless and keep the real features that beat them. The useless columns, called \emph{shadow features}, are random noise or copies of real features with their rows shuffled, which breaks any link to the target. After fitting a model, usually tree-based, any real feature whose importance falls below the largest shadow importance is dropped. Boruta \citep{kursa2010boruta} turns this into a procedure that repeats the comparison over many rounds and applies a binomial test to each feature's win count; BorutaPy \citep{homola2015borutapy} is its common Python implementation and BorutaSHAP \citep{keany2020borutashap} replaces impurity importance with SHAP values. Many teams also use the one-shot version directly on the gain importance of a gradient-boosted model. The appeal is that the shadows calibrate the threshold automatically: whatever importance pure noise can reach is, by construction, the bar to clear.

Credit default, fraud, medical diagnosis and churn all have rare positive classes, and in this setting the recipe drops features that do predict the target. This paper explains when and why, and what to do instead. The quantity that matters is the number of minority-class samples $\nmin$, combined with a feature's effect size. For a feature whose mean differs by $\delta$ within-class standard deviations between the classes, the combination is the \emph{scaling variable} $s = \nmin\delta^2$; for features of any distribution, a rank version $\sauc$ computed from the feature's ROC area plays the same role. A credit model with 5\% defaults among 1{,}000 loans and a fraud model with 0.1\% fraud among 50{,}000 transactions both have $\nmin = 50$, and to first order face the same difficulty.

The mechanism is not that noise features gain importance as the prevalence falls. The shadow rule keeps an irrelevant feature with probability at most $1/(K+1)$ for $K$ shadows at any imbalance (\cref{prop:level}). What changes is the advantage of an informative feature, which at the first split grows with $s$ while the best split pure noise can find depends on $\nmin$ alone and grows only slowly with it. Imbalance therefore takes away power and leaves the false positive level unchanged.

\paragraph{Contributions.}
\begin{enumerate}
\item \emph{Theory for the first split and for the shadow comparison (\cref{sec:theory}).} In the limit of small prevalence, the root split of a gradient-boosted tree depends on the data only through $\nmin$ and the positives' placement among the negatives (\cref{thm:limit}); for small effects on a fixed threshold window, only through $s$ (\cref{thm:scaling}). On XGBoost's full threshold range we conjecture, with numerical evidence, a slow $\log\log\nmin$ drift (\cref{conj:drift}). The shadow rule's level is $1/(K+1)$ for any learner whose importances depend on ranks, provided columns are continuous; per-feature shadows extend this to discrete columns (\cref{prop:level}). A rank test uses the same counting process as the split, has power that depends only on $\sauc$ for any continuous feature (\cref{prop:mw}), and is more efficient than any fixed split for Gaussian-like features but not for heavy-tailed or skewed ones (\cref{cor:efficiency}). XGBoost 1.x and 2.x intercepts add a toll of order $\lambda/\phat$ to every split (\cref{prop:toll}); 3.x does not.
\item \emph{Experiments that test each statement against pre-stated acceptance criteria (\cref{sec:experiments}, \cref{app:agreement}),} including the residual effect of the prevalence at fixed $\nmin$, which is real at weak signal.
\item \emph{A comparison of selection methods at both error rates (\cref{sec:methods}),} with SHAP and total-gain importances, BorutaPy, out-of-fold permutation importance, model-X knockoffs, $\ell_1$ and stability selection, and a rank-test screen, on synthetic data, on eight departures from the Gaussian model, and on five real tables with known null columns.
\item \emph{Guidance and a tested, scikit-learn compatible package (\cref{sec:guidance}),} \pkg{}, with an imbalance diagnostic, a rank-test screen and a Boruta-style selector. Every number and figure is produced by scripts in the repository with fixed seeds.
\end{enumerate}

\paragraph{What is classical and what is new.} Several ingredients are classical, and we credit them where they appear: the first-tree gain is Rao's score statistic, which for a binary split is Pearson's $X^2$ (\cref{prop:chi2}); the null best split is a maximally selected chi-square statistic, and its limit under local alternatives is the Brownian-bridge limit of \citet{miller1982maximally} and \citet{lausen1992maximally} (\cref{thm:scaling}); the level of the shadow rule is the Monte Carlo test argument of \citet{barnard1963discussion} and \citet{besag1989generalized} (\cref{prop:level}); and the efficiency of a fixed split relative to the rank test is the Pitman efficiency of the median test \citep{mood1954asymptotic}. New are the limit with the minority count fixed and the prevalence going to zero, which identifies $\nmin$ rather than $n$ as the sample size of the comparison; the scaling variables $s$ and $\sauc$ as the quantities that set the shadow rule's power; the intercept toll; per-feature shadows for discrete columns; and the assembly of these results into an account of shadow-based selection under imbalance, tested quantitatively.

\subsection{Related work}\label{sec:related}

\paragraph{Shadow, probe and knockoff methods.} Judging features against artificial ones that are irrelevant by construction predates Boruta: random probes \citep{stoppiglia2003ranking}, pseudo-variables that calibrate a selection threshold \citep{wu2007controlling}, artificial contrasts for tree ensembles \citep{tuv2009feature}, and pseudo-data that correct the bias of Gini importance \citep{sandri2008bias}. Boruta \citep{kursa2010boruta,kursa2011all} introduced permuted shadow copies with a random forest and a repeated binomial test, aimed at \emph{all-relevant} selection in the sense of \citet{nilsson2007consistent}; \citet{rudnicki2015all} review all-relevant methods. Benchmarks of Boruta against other random-forest selection methods, such as Vita \citep{janitza2018computationally}, VSURF \citep{genuer2010variable} and the method of \citet{hapfelmeier2013new}, include \citet{kursa2014robustness}, \citet{degenhardt2019evaluation} and \citet{speiser2019comparison}; none of them varies the minority count. Model-X knockoffs \citep{candes2018panning,barber2015controlling} compare each feature with its own synthetic copy and control the false discovery rate, so they do not pay the maximum-over-$K$ price of \cref{cor:shadows}; their power still depends on the importance statistic. Stability selection \citep{meinshausen2010stability,shah2013variable} controls false selections by subsampling a sparse model.

\paragraph{Feature importance in trees.} Impurity importances are biased toward features with many split points \citep{strobl2007bias} and can be analysed exactly for totally randomised trees \citep{louppe2013understanding}; debiased and out-of-sample versions exist \citep{li2019debiased,adler2022feature,zhang2023unbiased,scornet2023trees}. Permutation importance on out-of-bag or held-out data avoids the training-set bias \citep{breiman2001random,strobl2007bias}, can be made conditional \citep{strobl2008conditional}, and can be turned into p-values against a permutation null \citep{altmann2010permutation}. TreeSHAP gives consistent attributions, unlike the gain importance \citep{lundberg2017unified,lundberg2020local}. Split-selection bias motivates selecting splits by p-values rather than raw gain \citep{loh1997split,hothorn2006unbiased,strobl2007unbiased,nembrini2018revival}. Closest to our setting, \citet{janitza2013auc} show that error-rate permutation importance fails with a rare class and propose an AUC-based version, and balanced random forests keep every positive in every tree \citep{chen2004using}. These works address imbalance through the importance measure; we analyse how the minority count limits any shadow comparison.

\paragraph{Maximally selected statistics and rank tests.} The best split of a null feature is a maximally selected chi-square statistic \citep{miller1982maximally,lausen1992maximally,hothorn2003exact,hothorn2008generalized}. The placement values that represent the positives in our limit are standard in ROC analysis \citep{hanley1982meaning,delong1988comparing,pepe2004placement}. The loss from splitting a continuous variable at one threshold is the classical efficiency of grouping \citep{cox1957note} and of the median test relative to the Wilcoxon test \citep{mood1954asymptotic,mann1947test,hodges1956efficiency}.

\paragraph{Imbalanced learning.} Resampling, such as SMOTE \citep{chawla2002smote}, cost-sensitive learning and ensembles are surveyed by \citet{fernandez2018learning}. That literature targets predictive performance rather than the reliability of feature selection.
